% !TEX root = ../rl_book.tex
\chapter{Value-Based Methods}

\section*{Value Iteration and Its Connection to Q-Learning}

A fundamental dynamic programming method in reinforcement learning is
\emph{value iteration}, which computes the optimal value function $V^*(s)$
through repeated application of the Bellman optimality operator:
\[
V_{k+1}(s) = \max_a \Big[ R(s,a) + \gamma \, \mathbb{E}_{s'\sim P(\cdot|s,a)} 
               \big[ V_k(s') \big] \Big].
\]
By iteratively updating all states until convergence, value iteration recovers
$V^*$, and the corresponding optimal policy is obtained by choosing actions
that maximize the one-step lookahead.

Equivalently, one can update action-values via
\[
Q_{k+1}(s,a) = R(s,a) + \gamma \, \mathbb{E}_{s'\sim P(\cdot|s,a)} 
               \big[ \max_{a'} Q_k(s',a') \big],
\]
which directly yields $Q^*(s,a)$ without explicitly maintaining $V$.
This update requires knowledge of the transition dynamics $P(s'|s,a)$ and
reward function $R(s,a)$, and is thus a \emph{model-based} procedure.

\medskip
\noindent
\textbf{Connection to Q-Learning.}  
Q-learning can be understood as a \emph{model-free stochastic approximation} to
value iteration. Instead of requiring expectations with respect to the
transition model, Q-learning uses sampled transitions $(s,a,r,s')$ collected by
interacting with the environment:
\[
Q(s,a) \;\leftarrow\; Q(s,a) + \alpha \Big( r + \gamma \max_{a'} Q(s',a') - Q(s,a) \Big).
\]
This update mirrors the Bellman optimality operator but replaces the
expectation with a single observed transition. As training proceeds, the
estimates $Q(s,a)$ converge toward the optimal action-value function $Q^*(s,a)$
under suitable conditions. Thus, Q-learning can be viewed as the
model-free counterpart of value iteration, extending its principles to
situations where the environment dynamics are unknown and must be learned
purely from interaction.


\section{Deep Q-Network (DQN)}

A landmark in deep reinforcement learning is the Deep Q-Network (DQN) algorithm
\cite{mnih2015human}, which combines Q-learning with deep neural networks.
DQN approximates the optimal action-value function $Q^*(s,a)$ by a parameterized
network $Q_\theta(s,a)$, trained to minimize the Bellman error. A key innovation
was the use of two stabilizing techniques:

\begin{itemize}
    \item \textbf{Experience replay}: transitions $(s,a,r,s')$ are stored in a buffer
    and sampled uniformly for updates, breaking correlations between sequential data.
    \item \textbf{Target networks}: a separate set of parameters $\theta^{-}$ is held
    fixed for computing Bellman targets, updated only periodically, reducing the risk
    of divergence during training.
\end{itemize}

The learned Q-function induces a policy by greedy action selection,
\[
a(s) = \arg\max_a Q_\theta(s,a).
\]

Applied to the Atari 2600 benchmark from raw pixels, DQN demonstrated that a
single agent architecture can learn to play many diverse games directly from
high-dimensional sensory input, achieving performance comparable to or
surpassing human experts on several titles. This result established DQN as a
cornerstone of deep RL and catalyzed much of the subsequent progress in the
field.


% Preamble (once)
% \usepackage{algorithm}
% \usepackage{algpseudocode}
% \usepackage{float}
% \usepackage{amsmath}

\begin{algorithm}[H]
\caption{DQN with Replay, Periodic Batch Updates, and Target Network}
\label{alg:dqn-periodic}
\begin{algorithmic}[1]
\State \textbf{Hyperparams:} replay capacity $N$, batch size $B$, discount $\gamma$, lr $\alpha$
\State \hspace{2.6em} warm-up steps $N_{\text{warm}}$, update period $U$, updates-per-update $K$, target period $C$
\State Initialize replay buffer $\mathcal{D}$ to capacity $N$
\State Initialize Q-network $Q_\theta$ (random $\theta$); target network $Q_{\bar\theta}$ with $\bar\theta \gets \theta$
\For{episode $=1,2,\dots$}
  \State Initialize state $s$ (optionally preprocessed feature $\phi(s)$)
  \While{episode not done}
    \State Select $a$ by $\epsilon$-greedy w.r.t.\ $Q_\theta(s,\cdot)$
    \State Execute $a$; observe $(r, s', d)$ where $d=1$ iff terminal
    \State Push $(s,a,r,s',d)$ into $\mathcal{D}$; set $s \gets s'$
    \State \textbf{(Periodic learning)} \If{$|\mathcal{D}|\ge B$ \textbf{and} $t \ge N_{\text{warm}}$ \textbf{and} $t \bmod U = 0$}
      \For{$k=1$ to $K$} \Comment{$K$ SGD steps per update (UTD)}
        \State Sample minibatch $\{(s_i,a_i,r_i,s'_i,d_i)\}_{i=1}^B \sim \mathcal{D}$
        \State Targets (MSE/Huber loss):
        \[
          y_i \;=\; r_i + \gamma(1-d_i)\,\max_{a'} Q_{\bar\theta}(s'_i,a') 
          \quad\text{\footnotesize (Double DQN: }Q_{\bar\theta}(s'_i,\arg\max_{a'}Q_\theta(s'_i,a'))\text{)}
        \]
        \State $L(\theta)=\tfrac{1}{B}\sum_{i=1}^B \ell\!\big(y_i - Q_\theta(s_i,a_i)\big)$; 
               \quad $\theta \leftarrow \theta - \alpha \nabla_\theta L(\theta)$
      \EndFor
    \EndIf
    \State \textbf{(Target sync)} \If{$t \bmod C = 0$} \quad $\bar\theta \leftarrow \theta$ \EndIf
    \State Increment global step $t \leftarrow t+1$
  \EndWhile
\EndFor
\end{algorithmic}
\end{algorithm}
